\subsection{An acyclic opposite corner on five P5 states}

\begin{proposition}[Proved: the five-state acyclic-opposite P5 gate]
\label{CP2:A-P5-five-state-acyclic-opposite}
Let $N$ be any finite monoid with rationally acyclic
classifying space. A left-$N$, right-$P_5$ biset on at most
five states cannot give a collage of type $\CP^2$.
\end{proposition}
\begin{proof}
Write $P=P_5$, $X$ for the collage space, and $E=B(N,S,P)$
for its bridge. A right $P$ fixed state is excluded by
Lemma~\ref{CP2:A-mixed-five-reductions}(1).
In its absence the rectangular-band normal form has only
regular two-point core orbits. A unique such orbit is
excluded, at any state count, by part (3) of the same lemma.
With five states there can be at most two core orbits.
If there are exactly two and no transient, the full-regular-core
lemma applies. It remains to consider the four core states
$(a,j)$, with $a\in\{0,1\}$ indexing the orbits and
$j\in\{0,1\}$ the column, and one transient $w$.

For each band row $i$ there is an orbit $a_i$ such that
$we_{ij}=(a_i,j)$. If $a_0\ne a_1$, this right act is
the regular right $P$ act: send $w$ to $1$ and
$(a_i,j)$ to $e_{ij}$. Every commuting left $N$ actor
is consequently left multiplication by an element of $P$,
and the whole left action is specified by a homomorphism
$N\to P$. The projection $B(N,P,P)\to BN$ is a homotopy
equivalence. Indeed its right-action bar is the contractible
regular action category, and the equivariant map to a point
induces an equivalence after the other bar realization.
The homotopy pushout $BP\leftarrow E\to BN$ is therefore
homotopy equivalent to $BP$, which is not a projective plane.

Now suppose $a_0=a_1$. Row interchange
$\sigma(e_{ij})=e_{1-i,j}$ fixes the right action on every
state and defines a category automorphism $F$ fixing $N,S$.
Let $U$ be the four-point core and $A=\{0,1\}$ its orbit set,
with the left $N$ action induced on that set. Tensor the
length-two left projective $P$ resolution of
Lemma~\ref{CP2:A-order-five-corner} with $\mathbb Q[S]$.
Its first differential is zero, because the two row actors
agree. Its second differential maps isomorphically from
the two copies of the column-zero core onto $\mathbb Q[U]$.
Quotienting that acyclic subcomplex gives the $N$ complex
\[
 \mathbb Q[A][0]\oplus T[1],
 \qquad T=\mathbb Q[S]/\mathbb Q[U],
\]
with zero differential. The explicit semilinear lift in
Lemma~\ref{CP2:A-mixed-five-reductions} acts by $+1$ on
the first summand and $-1$ on the second. Thus the core
bridge $E_U=B(N,U,P)\to E$ induces exactly the $+1$
summand in rational homology, including in degree three.

Let $X_U$ be the full-core subcollage. The natural
nullhomotopy of its projection $E_U\to BP$ from
Lemma~\ref{CP2:A-P5-full-regular-core} gives
\[
 X_U\simeq BP\vee
 \operatorname{hocofib}\bigl(B(N,A)\longrightarrow BN\bigr).
\]
To justify the second factor, the equivariant map from the
right core-action category to its orbit set has contractible
fibers, each the action category of the regular column act
$e_{00}P$. It is a homotopy equivalence before and after
the $N$ bar. The second wedge summand is therefore the
ordinary left $N$ act cone on the two-element set $A$.
Its degree-two products vanish by the unconditional binary
act theorem. Those on $BP$ vanish as well, so every
degree-two product on $X_U$ is zero.

Since $BN$ is rationally acyclic and $BP$ has no third or
fourth rational homology, the two pushout sequences identify
$H_4(X_U)$ and $H_4(X)$ with $H_3(E_U)$ and $H_3(E)$,
respectively. Hence the map $H_4(X_U)\to H_4(X)$ realizes
the entire $+1$ eigenspace of $F$.

If $X$ had type $\CP^2$, let $\alpha$ generate its rational
second cohomology. A zero restriction to $BP$ would make
its square zero by the relative-cup lemma, since its
restriction to $BN$ already vanishes. Otherwise the
$-1$ action on $H^2(BP)$ forces $F^*\alpha=-\alpha$.
Its square is $F$ invariant and therefore pairs to zero
with the $-1$ eigenspace of $H_4(X)$. It also pairs to zero
with the $+1$ eigenspace: that space is in the image of
$H_4(X_U)$, and the restriction of $\alpha^2$ to $X_U$
vanishes. Thus $\alpha^2=0$, a contradiction.
\end{proof}

\begin{corollary}[Proved: the opposite-order-eight P5 scope at eighteen]
\label{CP2:A-P5-eight-five-scope}
For the profile $(5,8,5)$ with positive source corner, it
suffices to test all $1205$ ordinary order-eight monoids with
at least two minimal rows and two minimal columns as the
opposite corner. Every one of those tables is required,
including the ones with zero second rational homology.
\end{corollary}
\begin{proof}
The positive five-element source is $P_5$. A one-column
opposite ideal is removed by exact projective-right orbit
comparison and the bounded-eleven act theorem. A one-row
ideal makes the opposite classifying space rationally
acyclic, so the preceding proposition applies. The remaining
ideals have trivial group by the order-eight budget and
shapes $2\times2$, $2\times3$, $3\times2$. Their complete
catalogs contain $1181+12+12=1205$ tables, exactly as in
Proposition~\ref{CP2:A-P5-eight-corner-scope}. This is an
input reduction; no new five-state finite search is claimed
by this corollary.
\end{proof}
